% Template:     Tesis LaTeX
% Documento:    Archivo principal
% Versión:      3.4.2 (07/02/2025)
% Codificación: UTF-8
%
% Autor: Pablo Pizarro R.
%        pablo@ppizarror.com
%
% Manual template: [https://latex.ppizarror.com/tesis]
% Licencia MIT:    [https://opensource.org/licenses/MIT]

% CREACIÓN DEL DOCUMENTO
\documentclass[
	english, % Idioma: spanish, english, etc.	
	letterpaper, oneside
]{book}

\usepackage{graphicx}
\usepackage{listings}
\usepackage{xcolor}
\usepackage{pdfpages}
\usepackage{tikz}

% Configuración de colores para Python
\definecolor{codegreen}{rgb}{0,0.6,0}
\definecolor{codegray}{rgb}{0.5,0.5,0.5}
\definecolor{codepurple}{rgb}{0.58,0,0.82}
\definecolor{backcolour}{rgb}{0.95,0.95,0.92}

\lstdefinestyle{pythonstyle}{
	backgroundcolor=\color{backcolour},
	commentstyle=\color{codegreen},
	keywordstyle=\color{magenta},
	numberstyle=\tiny\color{codegray},
	stringstyle=\color{codepurple},
	basicstyle=\ttfamily\footnotesize,
	breakatwhitespace=false,
	breaklines=true,
	captionpos=b,
	keepspaces=true,
	numbers=left,
	numbersep=5pt,
	showspaces=false,
	showstringspaces=false,
	showtabs=false,
	tabsize=4,
	language=Python
}

\lstset{style=pythonstyle}

\newcommand{\pythonfunc}[1]{\texttt{\textbf{#1}}}
\newcommand{\pythonarg}[1]{\texttt{\textit{#1}}}

\graphicspath{{./img/}{./pictures/}}

% INFORMACIÓN DEL DOCUMENTO
\def\documenttitle {Study of the role of gas in the formation of SMBH seed, through simulations of NSC with gas embedded}
\def\documentsubtitle {}
\def\degreetitle {
	Tesis para optar al grado de magíster en ciencias, mención astronomía
	\bigbreak\vspace{0.3cm}
}

\def\universityname {Universidad de Chile}
\def\universityfaculty {Facultad de Ciencias Físicas y Matemáticas}
\def\universitydepartment {Departamento de Astronomía}
\def\universitydepartmentimage {departamentos/uchile2}
\def\universitydepartmentimagecfg {height=3cm}
\def\universitylocation {Santiago de Chile}

% INTEGRANTES, PROFESORES Y FECHAS
\def\documentauthor {Boris Cuevas Gomez}
\def\documentdate {\the\year}

\def\portrait {
	\begin{center}
	\vspace{1.5cm} ~ \\
	\MakeUppercase{\textbf{\documenttitle}} ~ \\
	\vspace{1.5cm}
	\MakeUppercase{\degreetitle} ~ \\
	\vfill
	\begin{tabular}{c}
		\MakeUppercase{\textbf{\documentauthor}} \\ \\
		\vspace{1.0cm} \\
		PROFESOR GUÍA: \\
		ANDRÉS ESCALA \\
		\vspace{0.5cm} \\
		MIEMBROS DE LA COMISIÓN: \\
		BASTIÁN REINOSO \\
		RICARDO MUÑOZ \\
		WALTER MAX-MOERBECK
		\vspace{0.5cm} \\
		Este trabajo ha sido parcialmente financiado por: \\
		NOMBRE INSTITUCIÓN \\
		\vspace{0.5cm} \\
		\MakeUppercase{\universitylocation} \\
		\MakeUppercase{\documentdate}
	\end{tabular}
	\end{center}
}
\def\abstracttable {
	\begin{tabular}{l}
		RESUMEN DE LA TESIS PARA OPTAR \\
		AL GRADO DE MAGÍSTER EN CIENCIAS, \\
		MENCIÓN ASTRONOMÍA \\
		POR: \MakeUppercase{\documentauthor} \\
		FECHA: \MakeUppercase{\documentdate} \\
		PROF. GUÍA: ANDRÉS ESCALA
	\end{tabular}
}

% IMPORTACIÓN DEL TEMPLATE
\input{template}

\makeatletter
\def\@cite#1#2{(#1\if@tempswa, #2\fi)}
\def\@biblabel#1{}
\makeatother

% INICIO DE LAS PÁGINAS
\begin{document}

% PORTADA
\templatePortrait

% CONFIGURACIÓN DE PÁGINA Y ENCABEZADOS
\templatePagecfg

% RESUMEN O ABSTRACT
\begin{abstractd}
	Observations of quasars at high redshift impose strong constraints on the formation mechanisms of supermassive black holes (SMBHs). Consequently, several pathways have been proposed to explain the formation of such massive objects ($\geq 10^6~M_{\odot}$) in the early universe. Among them, runaway collisions in nuclear star clusters (NSCs) are a promising channel for SMBH seed formation. In this work, we present a set of simulations of gas-embedded nuclear star clusters to study the effect of gas on the formation of a supermassive black hole seed. We explore a scenario that quantifies the role of collisions and the interplay with relaxation processes. The methods include an SPH code for the gas and an N-body code for the stars, both from the AMUSE framework, coupled externally via the BRIDGE method. We show that for simulations in the intermediate regime, the black hole formation efficiency is $\gtrsim 20\%$, and in our most massive models, an IMBH of $\geq 10^4~M_{\odot}$ is formed. Our findings show that gas enhances the formation of a SMBH seed by increasing the mass of the final most massive object and shortening the core collapse timescale.
\end{abstractd}

% DEDICATORIA
\begin{dedicatory}
	Una frase de dedicatoria, \\
	pueden ser dos líneas. \\
	~ \\
	\textbf{Saludos}
\end{dedicatory}

% AGRADECIMIENTOS
\begin{acknowledgments}
	Acknowledgments
\end{acknowledgments}

% TABLA DE CONTENIDOS - ÍNDICE
\templateIndex

% CONFIGURACIONES FINALES
\templateFinalcfg

% ======================= INICIO DEL DOCUMENTO =======================

\input{example} % Ejemplo, se puede borrar

% FIN DEL DOCUMENTO
\end{document}